% Options for packages loaded elsewhere
\PassOptionsToPackage{unicode}{hyperref}
\PassOptionsToPackage{hyphens}{url}
\PassOptionsToPackage{dvipsnames,svgnames,x11names}{xcolor}
%
\documentclass[
  10pt,
]{article}
\usepackage{amsmath,amssymb}
\usepackage{iftex}
\ifPDFTeX
  \usepackage[T1]{fontenc}
  \usepackage[utf8]{inputenc}
  \usepackage{textcomp} % provide euro and other symbols
\else % if luatex or xetex
  \usepackage{unicode-math} % this also loads fontspec
  \defaultfontfeatures{Scale=MatchLowercase}
  \defaultfontfeatures[\rmfamily]{Ligatures=TeX,Scale=1}
\fi
\usepackage{lmodern}
\ifPDFTeX\else
  % xetex/luatex font selection
  \setmainfont[]{DejaVu Serif}
  \setsansfont[]{DejaVu Sans}
  \setmonofont[]{DejaVu Sans Mono}
  \setmathfont[]{Latin Modern Math}
\fi
% Use upquote if available, for straight quotes in verbatim environments
\IfFileExists{upquote.sty}{\usepackage{upquote}}{}
\IfFileExists{microtype.sty}{% use microtype if available
  \usepackage[]{microtype}
  \UseMicrotypeSet[protrusion]{basicmath} % disable protrusion for tt fonts
}{}
\makeatletter
\@ifundefined{KOMAClassName}{% if non-KOMA class
  \IfFileExists{parskip.sty}{%
    \usepackage{parskip}
  }{% else
    \setlength{\parindent}{0pt}
    \setlength{\parskip}{6pt plus 2pt minus 1pt}}
}{% if KOMA class
  \KOMAoptions{parskip=half}}
\makeatother
\usepackage{xcolor}
\usepackage[margin=2.0cm]{geometry}
\usepackage{longtable,booktabs,array}
\usepackage{calc} % for calculating minipage widths
% Correct order of tables after \paragraph or \subparagraph
\usepackage{etoolbox}
\makeatletter
\patchcmd\longtable{\par}{\if@noskipsec\mbox{}\fi\par}{}{}
\makeatother
% Allow footnotes in longtable head/foot
\IfFileExists{footnotehyper.sty}{\usepackage{footnotehyper}}{\usepackage{footnote}}
\makesavenoteenv{longtable}
\setlength{\emergencystretch}{3em} % prevent overfull lines
\providecommand{\tightlist}{%
  \setlength{\itemsep}{0pt}\setlength{\parskip}{0pt}}
\setcounter{secnumdepth}{-\maxdimen} % remove section numbering
% doc_header.tex -- LaTeX preamble additions used by build_docs.sh for every
% Markdown -> LaTeX -> PDF conversion (README and provenance documents of Revision_3).
\usepackage{etoolbox}
% longtable may not be loaded by pandoc when a document has no table
\usepackage{longtable}
% bold vectors: \bm -> \symbf under unicode-math (XeLaTeX), bm package otherwise
\makeatletter
\@ifpackageloaded{unicode-math}{\providecommand{\bm}[1]{\symbf{#1}}}{\usepackage{bm}}
\makeatother
% smaller type in long tables so that wide audit tables fit the text width;
% inside a table a line break is allowed after an underscore of plain text
% (pandoc writes it as \_), so that long identifiers wrap inside their column
% rows are set a little apart, so that symbols with indices do not touch
\AtBeginEnvironment{longtable}{\footnotesize\renewcommand{\arraystretch}{1.45}%
  \let\sptunderscore\_\def\_{\sptunderscore\allowbreak}}
\setlength{\LTcapwidth}{\textwidth}
% allow line breaks after slashes and hyphens inside table cells
\usepackage{xurl}
% code blocks: break over-long lines instead of running into the margin.
% Plain ``` blocks are typeset by pandoc as {verbatim}, highlighted ones
% (```bash) as {Highlighting}; both are re-declared as breakable fancyvrb
% environments.  The continuation line is marked with a small hooked arrow.
\usepackage{fvextra}
\RecustomVerbatimEnvironment{verbatim}{Verbatim}{breaklines,breakanywhere,fontsize=\small}
\makeatletter
\@ifundefined{Highlighting}{}{%
  \RecustomVerbatimEnvironment{Highlighting}{Verbatim}{breaklines,breakanywhere,fontsize=\small,commandchars=\\\{\}}}
\makeatother
% inline code (pandoc writes it as \texttt{...}): never hyphenate -- an inserted
% hyphen could be read as part of a file name -- but allow a line break after
% an underscore, slash, comma, colon or equals sign, so that long paths and
% column lists wrap inside the text width and inside table columns.
% Inside a formula (e.g. the name of a notebook symbol in an equation) the same
% command must first leave math mode, which \text of amsmath does.
\ExplSyntaxOn
\tl_new:N \l__spt_code_tl
\cs_new_protected:Npn \__spt_texttt:n #1
  {
    \tl_set:Nn \l__spt_code_tl {#1}
    \regex_replace_all:nnN { ([\/\,\:\=]) } { \1 \c{allowbreak} } \l__spt_code_tl
    \regex_replace_all:nnN { \c{_} } { \c{_} \c{allowbreak} } \l__spt_code_tl
    \group_begin: \ttfamily \hyphenchar \font = -1 \scan_stop: \l__spt_code_tl \group_end:
  }
\RenewDocumentCommand \texttt { m }
  {
    \mode_if_math:TF { \text { \__spt_texttt:n {#1} } } { \__spt_texttt:n {#1} }
  }
\ExplSyntaxOff
% A figure is written in the Markdown source as an image paragraph followed by
% a caption paragraph ("**Figure S1.** ..."), so that the caption is also
% visible on GitHub.  An image that starts a paragraph is centered.
\usepackage{graphicx}% pandoc loads it only for documents that contain an image
\NewCommandCopy\sptincludegraphics\includegraphics
\RenewDocumentCommand\includegraphics{O{} m}{%
  \ifvmode\begin{center}\sptincludegraphics[#1]{#2}\end{center}%
  \else\sptincludegraphics[#1]{#2}\fi}
% keep figures where they are referenced
\usepackage{float}
\floatplacement{figure}{H}
% running header with the document title
\usepackage{fancyhdr}
\pagestyle{fancy}
\fancyhf{}
\fancyhead[L]{\small PROJECT spacetimes, Revision\_3}
\fancyhead[R]{\small\thepage}
\renewcommand{\headrulewidth}{0.4pt}
% 16x16 matrices: tighter columns so that a full matrix fits the text width
\setlength{\arraycolsep}{3.2pt}
% The documents of this project contain long tagged symbols, file paths and
% identifiers that cannot be hyphenated.  A paragraph is therefore allowed to
% have looser lines before TeX gives up and lets a line stick out of the margin.
\AtBeginDocument{\tolerance=2000 \emergencystretch=5em \hbadness=2000}
\ifLuaTeX
  \usepackage{selnolig}  % disable illegal ligatures
\fi
\IfFileExists{bookmark.sty}{\usepackage{bookmark}}{\usepackage{hyperref}}
\IfFileExists{xurl.sty}{\usepackage{xurl}}{} % add URL line breaks if available
\urlstyle{same}
\hypersetup{
  pdftitle={The symbols of Revision\_3: registry, collisions resolved, machine checks},
  pdfauthor={Prepared for Patrick Lee Nash, Ph.D.},
  colorlinks=true,
  linkcolor={blue},
  filecolor={Maroon},
  citecolor={blue},
  urlcolor={blue},
  pdfcreator={LaTeX via pandoc}}

\title{The symbols of Revision\_3: registry, collisions resolved,
machine checks}
\usepackage{etoolbox}
\makeatletter
\providecommand{\subtitle}[1]{% add subtitle to \maketitle
  \apptocmd{\@title}{\par {\large #1 \par}}{}{}
}
\makeatother
\subtitle{Provenance document \texttt{P1} of PROJECT spacetimes,
Revision\_3}
\author{Prepared for Patrick Lee Nash, Ph.D.}
\date{November 16, 2026}

\begin{document}
\maketitle

\hypertarget{what-this-document-is}{%
\section{1. What this document is}\label{what-this-document-is}}

Hard Rule 5 of the author requires that no symbol of PROJECT spacetimes
denote two objects, that a MASTER provenance of the symbols exist, and
that every document use it. This document is the registry of the symbols
of the folder \texttt{Revision\_3}: it states how the symbols were
chosen (Section 2), lists every collision with a symbol of the rest of
the project that was found and how it was resolved (Section 3),
describes the six machine checks and prints their complete output
(Section 4), and gives the complete table of the symbols (Section 5),
which every document of the folder also contains. The machine-readable
registry is the file \texttt{provenance/symbols\_r3.csv}, written by
\texttt{code/scripts/write\_symbol\_registry.py}.

\hypertarget{how-the-symbols-were-chosen}{%
\section{2. How the symbols were
chosen}\label{how-the-symbols-were-chosen}}

\begin{enumerate}
\def\labelenumi{\arabic{enumi}.}
\tightlist
\item
  \textbf{Taken symbols.} Thirteen entries are taken from the two master
  registries of the project (the registry of the main project and the
  MASTER registry of the revision on the deflating spacetimes): the
  coordinates \(X=\{x0,\dots,x7\}\), the index letters, the number \(D\)
  of space dimensions, the index sets \(K\) and ranks \(r\), and the
  matrices of the 7+1 spacetime (\(\eta_{1671}\), \(\gamma^A_{1671}\),
  \(\Sigma^{AB}_{1671}\), \(\omega_{1671}\), \(\chi_{1671}\),
  \(P^L_{1671}\), \(P^R_{1671}\), \(C^\pm_{1671}\), \(\beta_{1671}\),
  \(\alpha^j_{1671}\)). Their meaning is copied from the master file,
  word for word, so that it cannot drift; where Revision\_3 uses only
  part of an entry, a note is appended (for instance: only the index
  letters \(A\) and \(C\) are used, because \(B_{U2A}\) is a gauge field
  here).
\item
  \textbf{New symbols carry the suffix U2}, the name of the new gauge
  group \(U(2)_{U2}\), in their subscript: \(\Delta_{U2}\),
  \(\mu_{U2}\), \(\Omega_{U2}\), \(G^P_{U2}\), \(m_{U2}\),
  \(\sigma_{U2}\), \(B_{U2A}\), and so on. A new symbol without the
  suffix is either a mark that no master registry uses (sans-serif
  matrices such as \(\mathsf G_K\), \(\mathsf E\), \(\mathsf K\); the
  index \(\iota\)), or the real and imaginary parts \(\Delta_1\),
  \(\Delta_2\) of the pair field, or the imaginary unit \(i\), the
  component indices \(l\), \(l'\) and the momentum \(\vec q\), \(q\),
  \(q^j\), which no master registry defines with another meaning in the
  documents of the project that Revision\_3 relies on.
\item
  \textbf{The electroweak theory}, quoted as the template, carries the
  suffix SM.
\item
  \textbf{Names in programs and data files.} The programs write symbols
  as plain names (for instance \texttt{Delta} for \(\Delta_{U2}\),
  \texttt{mu} for \(\mu_{U2}\)); every column of every data file is
  either a registered name (last column of the table of Section 5) or a
  descriptive word. All data files of the folder begin with
  \texttt{r3\_} and state their meaning in a header line
  \texttt{\#\ meaning:}, so a name can only be read as the symbol of
  Revision\_3.
\end{enumerate}

\hypertarget{the-collisions-that-were-found-and-how-each-was-resolved}{%
\section{3. The collisions that were found, and how each was
resolved}\label{the-collisions-that-were-found-and-how-each-was-resolved}}

The left column names the letter as a plain name (it is the symbol of
another part of the project, not of Revision\_3); the right column gives
what Revision\_3 writes instead.

\begin{longtable}[]{@{}
  >{\raggedright\arraybackslash}p{(\columnwidth - 4\tabcolsep) * \real{0.3333}}
  >{\raggedright\arraybackslash}p{(\columnwidth - 4\tabcolsep) * \real{0.3333}}
  >{\raggedright\arraybackslash}p{(\columnwidth - 4\tabcolsep) * \real{0.3333}}@{}}
\toprule\noalign{}
\begin{minipage}[b]{\linewidth}\raggedright
letter or symbol in the rest of the project
\end{minipage} & \begin{minipage}[b]{\linewidth}\raggedright
its meaning there
\end{minipage} & \begin{minipage}[b]{\linewidth}\raggedright
what Revision\_3 writes
\end{minipage} \\
\midrule\noalign{}
\endhead
\bottomrule\noalign{}
\endlastfoot
\texttt{lambda}, \texttt{lambda\_S}, \texttt{lambda\_P} & the
dimensionless pairing coupling and its parts from the radial and the
phase mode & \(\lambda^P_{U2}=G^P_{U2}/G^{\rm crit}_{U2}\),
\(\lambda^S_{U2}=4G^S_{U2}/G^{\rm crit}_{U2}\) \\
\texttt{G\_S}, \texttt{G\_V} & the couplings of the Kohn--Sham
description of matter & \(G^S_{U2}\), \(G^P_{U2}\) \\
\texttt{Lambda}, \texttt{alpha\_c} & a momentum cutoff of the ladder
equation & \(\Lambda_{U2}\) \\
\texttt{mu}, \texttt{k\_F}, \texttt{E\_F} & chemical potential, Fermi
momentum and energy of matter & \(\mu_{U2}\), \(q^F_{U2}\) \\
\texttt{Delta(k)}, \texttt{Delta\^{}a}, \texttt{Delta\_1671\^{}A} & gap
function; vector order parameter; vector pair operators &
\(\Delta_{U2}\), \(\Delta_1\), \(\Delta_2\) \\
\texttt{Omega\_mu} & the spin connection of the deflating spacetime &
\(\Omega_{U2}\), \(\Omega^{\rm cl}_{U2}\), \(\Omega^{\min}_{U2}\),
\(\Omega^{\rm sym}_{U2}\) \\
\texttt{Theta\_K} (with a phase \texttt{zeta\_r}) & the 256 products of
Dirac matrices of the self-interaction & \(\mathsf G_K\), the ordered
products without a phase \\
\texttt{Phi}, \texttt{Phi(p)} & a state of a pair, an antisymmetric
matrix & \(\Phi^\pm_{U2}\), the pair operators \\
\texttt{tau}, \texttt{tau\_d} & optical depth; deflation times &
\(\tau^1_{U2}\), \(\tau^2_{U2}\), \(\tau^3_{U2}\), the Pauli matrices \\
index \texttt{B}; \texttt{B} (bag constant) & an index with the values 0
to 7; an energy density & only \(A\) and \(C\) as indices; \(B_{U2A}\)
the \(U(1)_{U2}\) gauge field \\
\texttt{b\_5}, \texttt{b\_6}, \texttt{b\_7}; \texttt{a\_s},
\texttt{a\_sc} & scale factors; other constants & \(a^{(7)}_{U2}\),
\(a^{(8)}_{U2}\) \\
\texttt{m} (with \texttt{beta}, \texttt{alpha\^{}j}, \texttt{chi},
\texttt{h}) & the mass of the spacetime at hand & \(m_{U2}\), the bare
mass of the doublet \\
\texttt{M\_1671}, \texttt{M\_s}, \texttt{mathsf\ M} & matrices and
masses of other parts & \(M_{U2}\), the effective mass \\
\texttt{x} & the ratio of masses \texttt{m*/m} & only the coordinate
names \(x0,\dots,x7\), \(x^A\), \(x^j\) and \(\vec x\) \\
\texttt{w} & a constant spinor of the exact solutions on the deflating
spacetime & \(w_{U2}\), a frequency; eigenvectors are
\(\mathsf e^\pm_{U2}\) \\
\texttt{s}, \texttt{t} & indices of the coordinates \(x5\), \(x6\),
\(x7\) & \(s_{U2}=M_{U2}^2+\lvert\Delta_{U2}\rvert^2\); no \texttt{t} \\
\texttt{a}, \texttt{b} & indices of the four extra axes & not used;
spinor components are indexed by \(l\), \(l'\) \\
\texttt{l} & the summation index of the closed form of the Fierz matrix
& \(l\), \(l'\) index spinor components; the Fierz matrix does not occur
in Revision\_3 \\
\texttt{K} & an index set (taken) & an index set only; a matrix of a
bilinear form is \(\mathsf K\) \\
\texttt{S}, \texttt{S\_0} & the propagator of the CJT effective action &
``spin zero'' is written in words \\
\texttt{i} & an index of space-like coordinates & the imaginary unit;
space-like indices are \(j\), \(j'\) \\
\texttt{n}, \texttt{n\_S} & number and scalar density of matter &
\(n_{U2}\) \\
\texttt{N} & the density of dark radiation & \(N_{U2}\), the number
operator \\
\texttt{E\_k}, \texttt{xi\_k} & quasi-particle energy, kinetic energy &
\(E^\pm_{U2}\), \(\xi^\pm_{U2}\) \\
\texttt{epsilon(k\_F,x)}, \texttt{epsilon\_kin} & energy densities &
\(\varepsilon_{U2}\), the single-particle energy \\
\texttt{kappa(k)}, \texttt{kappa\_8}, \texttt{kappa\_SI} & a pair
correlation; constants & \(\kappa_{U2}\), \(\kappa^{(4)}_{U2}\), the
Euclidean momentum \\
\texttt{sigma(p)}, \texttt{sigma\_v}, \texttt{sigma\_8} & a self-energy;
a velocity dispersion; a clustering amplitude & \(\sigma_{U2}\), the
amplitude mode \\
\texttt{theta\_s} & an angle of cosmology & \(\theta_{U2}\), the phase
of the pair field \\
\texttt{f\_+}, \texttt{f\_D}, \texttt{f\_g} & functions and constants &
\(f^s_{U2}\), \(f^t_{U2}\), \(f_{U2}\) \\
\texttt{g\_r}, \texttt{g\_VS} & couplings of the self-interaction and of
the Vorton-Scalar & \(g_{U2}\), \(g'_{U2}\) \\
\texttt{y\_VS} & a charge of the Vorton-Scalar & \(y_{U2}\) \\
\texttt{z\_*}, \texttt{z\_t} & redshifts & \(z_{U2}\) \\
\texttt{u\_d} & a scaled wave number & \(u^\pm_{U2}\), \(v^\pm_{U2}\) \\
\texttt{Pi\_1671} & a projector & \(\Pi_{U2}\), the polarization \\
\texttt{R} & the curvature scalar & \(\mathsf R_{U2}\), the matrix of
small oscillations \\
\texttt{V\_d}, \texttt{V\_0}, \texttt{V\_s(r)} & a volume; potentials &
\(V^{\rm C}_{U2}\), \(V^{\rm pair}_{U2}\), \(\mathsf V^1_{U2}\) to
\(\mathsf V^3_{U2}\), \(\mathrm{Vol}_{U2}\) \\
\texttt{varrho\_c} & an abbreviation of a ladder equation &
\(\varrho_{U2}\), a distance \\
\texttt{phi\_1671\^{}L}, \texttt{phi\_bos} & chiral parts; a boson field
& \(\phi^{\rm Y}_{U2}\), \(\phi_{\rm SM}\) \\
\texttt{psi\_1671} & the free field dirac1671 & \(\psi^1_{U2}\),
\(\psi^2_{U2}\), \(\Psi_{U2}\) \\
\texttt{L}, \texttt{R} & chirality labels & kept as labels of
\(P^L_{1671}\), \(P^R_{1671}\); \(L_{U2}\) is a function of (R8) \\
\texttt{J} & an index set & \(J^A_{U2}\), \(\boldsymbol J^A_{U2}\)
(currents), \(\mathcal J^{(\iota)}_{U2}\) (integrals) \\
\texttt{c\_s} & a sound speed & \(c^\sigma_{U2}\) \\
\texttt{Q(M)} & a function of a mass & \(\vec Q_{U2}\) \\
\texttt{Y\_He} & the helium fraction & \(Y_{\rm SM}\) \\
\end{longtable}

\hypertarget{the-machine-checks-of-the-symbols}{%
\section{4. The machine checks of the
symbols}\label{the-machine-checks-of-the-symbols}}

\texttt{code/scripts/check\_symbols.py} verifies:

\begin{itemize}
\tightlist
\item
  \textbf{R3-S1}: no symbol text is registered twice, and no name used
  in a data file belongs to two entries;
\item
  \textbf{R3-S2}: every taken entry carries exactly the meaning of the
  master registry, and no new symbol text is a symbol text of either
  master registry;
\item
  \textbf{R3-S3}: every column of every data file is a registered name
  or a descriptive word, and every data file states its meaning;
\item
  \textbf{R3-S4}: every document of the folder, and the manuscript,
  contains the complete table of the symbols;
\item
  \textbf{R3-S5}: in every formula of the documents, \(\lambda^P_{U2}\),
  \(\lambda^S_{U2}\), \(G^P_{U2}\), \(G^S_{U2}\), \(G^{\rm crit}_{U2}\),
  \(\mu_{U2}\), \(\Lambda_{U2}\), \(\Omega_{U2}\) and \(\Delta_{U2}\)
  carry their suffix;
\item
  \textbf{R3-S6}: in every formula of the documents and of the new
  registry entries, no letter that a master registry defines with
  another meaning stands without the suffix of its Revision\_3 meaning;
  the only letters allowed without a suffix are the taken indices and
  sets (\(A\), \(C\), \(D\), \(K\), \(X\), \(j\), \(r\)), the new plain
  symbols (\(i\), \(l\), \(q\)) and the differential \(d\); the letter x
  occurs only in the coordinate names.
\end{itemize}

The complete output of the last run:

\begin{verbatim}
PROJECT spacetimes, Revision_3: checks of the symbols (registry provenance/symbols_r3.csv)
6 of 6 checks passed

[PASS] R3-S1  every symbol text is registered once (90 entries) and every name used in the data files of Revision_3 belongs to one entry (88 names)
[PASS] R3-S2  the 13 taken entries carry exactly the meaning of the MASTER registry (113 entries of the main project, 229 of the revision on the deflating spacetimes compared), and none of the 77 new symbols is a symbol of either master registry
[PASS] R3-S3  every column of the 18 data files data/r3_*.csv is a registered name or a descriptive word (88 different columns), and every file states its meaning; 9 column names also occur as names in a master registry or as symbols of the author's notebooks (A, M, m, mu, n, name, q, rank, x) - they cannot be confused, because every data file of Revision_3 has the prefix r3_ and a line '# meaning:', and the name means in it the registered symbol of Revision_3
[PASS] R3-S4  each of the 4 documents of Revision_3 (and the manuscript) contains the complete table of the 90 registered symbols, so that every symbol it uses is defined in the document itself
[PASS] R3-S5  in the 1554 formulas of the documents, lambda^P, lambda^S, G^P, G^S, G^crit, mu, Lambda, Omega and Delta always carry the suffix U2 (Delta also the index 1 or 2 of its real and imaginary parts)
[PASS] R3-S6  in the 1910 formulas of the documents and of the new registry entries, no letter that a master registry defines with another meaning stands bare: every letter except the taken indices and sets (A, C, D, K, X, j, r), the new plain symbols (i, l, q) and the differential d carries the suffix of its Revision_3 meaning, x occurs only in the coordinate names, and the Greek letters that are master symbols carry their suffix
\end{verbatim}

\hypertarget{the-complete-table-of-the-symbols}{%
\section{5. The complete table of the
symbols}\label{the-complete-table-of-the-symbols}}

``taken'': a symbol of the master registries, used with the same meaning
(the master meaning is copied, followed by a note on its use in
Revision\_3). ``new'': introduced in Revision\_3.

\textbf{1 Spacetime, coordinates and indices}

\begin{longtable}[]{@{}
  >{\raggedright\arraybackslash}p{(\columnwidth - 6\tabcolsep) * \real{0.2200}}
  >{\raggedright\arraybackslash}p{(\columnwidth - 6\tabcolsep) * \real{0.5000}}
  >{\raggedright\arraybackslash}p{(\columnwidth - 6\tabcolsep) * \real{0.0800}}
  >{\raggedright\arraybackslash}p{(\columnwidth - 6\tabcolsep) * \real{0.2000}}@{}}
\toprule\noalign{}
\begin{minipage}[b]{\linewidth}\raggedright
symbol
\end{minipage} & \begin{minipage}[b]{\linewidth}\raggedright
meaning
\end{minipage} & \begin{minipage}[b]{\linewidth}\raggedright
origin
\end{minipage} & \begin{minipage}[b]{\linewidth}\raggedright
names in programs and data files
\end{minipage} \\
\midrule\noalign{}
\endhead
\bottomrule\noalign{}
\endlastfoot
\(X=\{x0,\dots,x7\}\) & the eight coordinates of both spacetimes, never
renamed: \(x0\) hidden space-like; \(x1,x2,x3\) ordinary space; \(x4\)
ordinary time; \(x5,x6,x7\) additional coordinates (time-like in 4+4,
space-like in 7+1) & taken & \texttt{x0} \texttt{x1} \texttt{x2}
\texttt{x3} \texttt{x4} \texttt{x5} \texttt{x6} \texttt{x7} \\
\(A,B,C,D\) & index with the values \(0,\dots,7\); always labels the
coordinate \(xA\). In Revision\_3 only the letters \(A\) and \(C\) are
used as indices; \(B_{U2A}\) is the \(U(1)_{U2}\) gauge field. & taken
& \\
\(i,j\) & index of space-like coordinates; values 1, 2, 3 (ordinary
space) unless a sum is said to run over all space-like coordinates of
the spacetime. In Revision\_3 only \(j\) (and a second such index
\(j'\)) is used as an index, and sums over \(j\) run over all seven
space-like coordinates \(0,1,2,3,5,6,7\); the letter \(i\) is the
imaginary unit. & taken & \\
\(D\) & the number of space dimensions in which the particles of matter
have momentum: 7 (all space coordinates of 7+1), 4 (fields that do not
depend on \(x5,x6,x7\)) or 3 (fields that depend on ordinary space
only). \(V_D=\pi^{D/2}/\Gamma_{\rm E}(D/2+1)\) is the volume and
\(A_D=D\,V_D\) the area of the unit ball in \(D\) dimensions
(\(\Gamma_{\rm E}\): Euler's gamma function). In Revision\_3 \(D=7\)
(\(A_7=16\pi^3/15\)), and \(A_8=\pi^4/3\) is the area of the unit sphere
of the Euclidean 8-momentum. & taken & \texttt{D} \\
\(K\), \(J\) (index sets); \(r\), \(r'\) (ranks) & \(K\), \(J\): a
strictly increasing list \(A_1<\dots<A_r\) of indices out of
\(0,\dots,7\) (there are 256); \(r=\lvert K\rvert\) is its rank,
\(0\le r\le8\); \(r'\) is a second rank and \(l\) the summation index in
the closed form of the Fierz matrix. In the data files an index set is
given by its bit mask (\texttt{mask}: the sum of \(2^A\) over its
indices). In Revision\_3 only \(K\) and \(r\) are used, for the ordered
products \(\mathsf G_K\); the Fierz matrix does not occur. & taken &
\texttt{mask} \texttt{rank} \texttt{J} \\
\(\vec x\) & the seven space-like coordinates \((x0,x1,x2,x3,x5,x6,x7)\)
as one vector; \(\vec q\cdot\vec x=\sum_jq^jx^j\) & new & \\
\(l\), \(l'\) & indices of the components of a spinor (values
\(1,\dots,16\)) or of the doublet \(\Psi_{U2}\) (values \(1,\dots,32\));
the summation index of the closed form of the Fierz matrix, which the
master registry also calls \(l\), does not occur in Revision\_3 & new
& \\
\(\delta^{jj'}\) & the Kronecker symbol: 1 if \(j=j'\), 0 otherwise &
new & \\
\(\varrho_{U2}\) & the distance between two static sources in the
seven-dimensional space & new & \\
\end{longtable}

\textbf{2 Matrices of the 7+1 spacetime}

\begin{longtable}[]{@{}
  >{\raggedright\arraybackslash}p{(\columnwidth - 6\tabcolsep) * \real{0.2200}}
  >{\raggedright\arraybackslash}p{(\columnwidth - 6\tabcolsep) * \real{0.5000}}
  >{\raggedright\arraybackslash}p{(\columnwidth - 6\tabcolsep) * \real{0.0800}}
  >{\raggedright\arraybackslash}p{(\columnwidth - 6\tabcolsep) * \real{0.2000}}@{}}
\toprule\noalign{}
\begin{minipage}[b]{\linewidth}\raggedright
symbol
\end{minipage} & \begin{minipage}[b]{\linewidth}\raggedright
meaning
\end{minipage} & \begin{minipage}[b]{\linewidth}\raggedright
origin
\end{minipage} & \begin{minipage}[b]{\linewidth}\raggedright
names in programs and data files
\end{minipage} \\
\midrule\noalign{}
\endhead
\bottomrule\noalign{}
\endlastfoot
\(\eta_{1671}\) & metric of the 7+1 spacetime,
\(\mathrm{diag}(-1,-1,-1,-1,+1,-1,-1,-1)\): time-like \(=+1\), time at
index 4 & taken & \texttt{eta1671} \\
\(\gamma_{1671}^A\) & the eight real \(16\times16\) Dirac matrices of
the 7+1 spacetime (contravariant index) & taken & \texttt{gamma1671} \\
\(\Sigma_{1671}^{AB}\) & \(\tfrac14[\gamma_{1671}^A,\gamma_{1671}^B]\),
the scaled commutators; they span \(\mathfrak{so}(7,1)\) & taken &
\texttt{Sigma1671} \\
\(\omega_{1671}\) & the product
\(\gamma_{1671}^0\cdots\gamma_{1671}^7\); real, \(\omega_{1671}^2=-1\) &
taken & \texttt{omega1671} \\
\(\chi_{1671}\) & \(i\,\omega_{1671}\), the (complex) chirality matrix
of 7+1 & taken & \texttt{chi1671} \\
\(P_{1671}^{L}\), \(P_{1671}^{R}\) & the complex chiral projectors
\(\tfrac12(1\mp\chi_{1671})\) & taken & \texttt{PL1671}
\texttt{PR1671} \\
\(C_{1671}^{+}\), \(C_{1671}^{-}\) & \(\gamma_{1671}^4\) and
\(\gamma_{1671}^4\omega_{1671}\);
\(C^{\pm}\gamma^A(C^{\pm})^{-1}=\pm(\gamma^A)^T\) & taken &
\texttt{Cp1671} \texttt{Cm1671} \\
\(\beta_{1671}\), \(\alpha_{1671}^j\) & the real symmetric matrices
\(\gamma_{1671}^4\) and \(\gamma_{1671}^4\gamma_{1671}^j\) (seven space
indices \(j\)) of the Hamiltonian & taken & \texttt{beta1671}
\texttt{alpha1671} \\
\(\mathsf G_K\) & the ordered product
\(\gamma^{A_1}_{1671}\gamma^{A_2}_{1671}\cdots\gamma^{A_r}_{1671}\) over
an index set \(K=(A_1<\dots<A_r)\) (\(\mathsf G_K=\mathsf 1_{16}\) for
the empty set); the 256 matrices \(\mathsf G_K\) are a basis of the real
\(16\times16\) matrices & new & \\
\(\mathsf 1_2\), \(\mathsf X\), \(\mathsf Z\), \(\mathsf E\) & the real
\(2\times2\) matrices \(\mathsf 1_2=\mathrm{diag}(1,1)\),
\(\mathsf X=\begin{pmatrix}0&1\\1&0\end{pmatrix}\),
\(\mathsf Z=\begin{pmatrix}1&0\\0&-1\end{pmatrix}\),
\(\mathsf E=\begin{pmatrix}0&1\\-1&0\end{pmatrix}\); each
\(\gamma^A_{1671}\) is \(\pm\) a Kronecker product of four of them, and
on the doublet index \(\mathsf E=i\tau^2_{U2}\) is the antisymmetric
matrix of the spin-zero pair (the columns f1 to f4 of
r3\_gamma\_kron.csv hold the letters 1, X, Z, E) & new & \texttt{1}
\texttt{X} \texttt{Z} \texttt{E} \\
\(\mathsf 1_{16}\), \(\mathsf 1_{32}\) & the unit matrices of the 16
spinor components and of the 32 components of the doublet (or of the
Nambu spinor) & new & \\
\(\mathsf K\) & a real \(16\times16\) matrix that defines the bilinear
form \(\psi^{1T}_{U2}\mathsf K\psi^2_{U2}\) of two spinors (Sec.~on the
pairs) & new & \\
\end{longtable}

\textbf{3 The new gauge sector: fields, group, covariant derivative}

\begin{longtable}[]{@{}
  >{\raggedright\arraybackslash}p{(\columnwidth - 6\tabcolsep) * \real{0.2200}}
  >{\raggedright\arraybackslash}p{(\columnwidth - 6\tabcolsep) * \real{0.5000}}
  >{\raggedright\arraybackslash}p{(\columnwidth - 6\tabcolsep) * \real{0.0800}}
  >{\raggedright\arraybackslash}p{(\columnwidth - 6\tabcolsep) * \real{0.2000}}@{}}
\toprule\noalign{}
\begin{minipage}[b]{\linewidth}\raggedright
symbol
\end{minipage} & \begin{minipage}[b]{\linewidth}\raggedright
meaning
\end{minipage} & \begin{minipage}[b]{\linewidth}\raggedright
origin
\end{minipage} & \begin{minipage}[b]{\linewidth}\raggedright
names in programs and data files
\end{minipage} \\
\midrule\noalign{}
\endhead
\bottomrule\noalign{}
\endlastfoot
\(\Psi_{U2}\); \(\psi^1_{U2}\), \(\psi^2_{U2}\) & the doublet of the new
gauge theory: two quantized 16-component Dirac fields \(\psi^1_{U2}\),
\(\psi^2_{U2}\) on the 7+1 spacetime whose components are complex
Grassmann operators; \(\Psi_{U2}=(\psi^1_{U2},\psi^2_{U2})\), 32
components, doublet index outside & new & \\
\(\bar\Psi_{U2}\); \(\bar\psi^1_{U2}\), \(\bar\psi^2_{U2}\) & the Dirac
adjoint taken member by member,
\(\bar\Psi_{U2}=(\bar\psi^1_{U2},\bar\psi^2_{U2})=(\psi^{1\dagger}_{U2}\beta_{1671},\ \psi^{2\dagger}_{U2}\beta_{1671})\)
& new & \\
\(\tilde\Psi_{U2}\) & the conjugate doublet
\(\mathsf E\,\Psi^*_{U2}=(\psi^{1*}_{U2},\psi^{2*}_{U2})\) multiplied by
\(\mathsf E\), i.e.~\((\psi^{2*}_{U2},-\psi^{1*}_{U2})\); since the
\(\gamma^A_{1671}\) are real, \(\Psi^*_{U2}\) transforms under Spin(7,1)
like \(\Psi_{U2}\), and \(\tilde\Psi_{U2}\) transforms under
\(SU(2)_{U2}\) like \(\Psi_{U2}\), with \(U(1)_{U2}\) charge \(-y_{U2}\)
& new & \\
\(\tau^1_{U2}\), \(\tau^2_{U2}\), \(\tau^3_{U2}\);
\(\boldsymbol\tau_{U2}\) & the Pauli matrices
\(\begin{pmatrix}0&1\\1&0\end{pmatrix}\),
\(\begin{pmatrix}0&-i\\i&0\end{pmatrix}\),
\(\begin{pmatrix}1&0\\0&-1\end{pmatrix}\) acting on the doublet index;
\(\boldsymbol\tau_{U2}=(\tau^1_{U2},\tau^2_{U2},\tau^3_{U2})\); the
three matrices \(\tfrac12\tau^1_{U2}\), \(\tfrac12\tau^2_{U2}\),
\(\tfrac12\tau^3_{U2}\) generate \(SU(2)_{U2}\);
\(\boldsymbol\tau^{(1)}_{U2}\), \(\boldsymbol\tau^{(2)}_{U2}\): the same
matrices acting on the first and on the second of two static sources &
new & \\
\(U(2)_{U2}\); \(SU(2)_{U2}\), \(U(1)_{U2}\) & the new gauge group
\(U(2)_{U2}=(SU(2)_{U2}\times U(1)_{U2})/\mathbb Z_2\) acting on the
doublet index: \(SU(2)_{U2}\) by unitary \(2\times2\) matrices of
determinant 1, \(U(1)_{U2}\) by the phases
\(\exp(i y_{U2}\Theta^0_{U2})\); it commutes with Spin(7,1) & new & \\
\(y_{U2}\) & the \(U(1)_{U2}\) charge of \(\Psi_{U2}\),
\(y_{U2}=\tfrac12\); the pairs \(\Phi^\pm_{U2}\) carry \(2y_{U2}=1\) &
new & \texttt{Y} \\
\(g_{U2}\), \(g'_{U2}\) & the gauge couplings of \(SU(2)_{U2}\) and
\(U(1)_{U2}\); on the 7+1 spacetime of dimension mass\(^{-2}\) & new &
\texttt{G} \texttt{GP} \\
\(\boldsymbol W_{U2A}=(W^1_{U2A},W^2_{U2A},W^3_{U2A})\) & the
\(SU(2)_{U2}\) gauge field, a triple of real vector fields & new &
\texttt{w} \\
\(B_{U2A}\) & the \(U(1)_{U2}\) gauge field, a real vector field & new &
\texttt{b} \\
\(\boldsymbol W_{U2AC}\), \(B_{U2AC}\) & the field strengths
\(\boldsymbol W_{U2AC}=\partial_A\boldsymbol W_{U2C}-\partial_C\boldsymbol W_{U2A}-g_{U2}\,\boldsymbol W_{U2A}\times\boldsymbol W_{U2C}\)
(\(\times\) the cross product of triples) and
\(B_{U2AC}=\partial_AB_{U2C}-\partial_CB_{U2A}\) & new & \\
\(D_{U2A}\) & the new covariant derivative
\(D_{U2A}\Psi_{U2}=(\partial_A+\tfrac i2 g_{U2}\boldsymbol\tau_{U2}\cdot\boldsymbol W_{U2A}+i g'_{U2}y_{U2}B_{U2A})\Psi_{U2}\);
it acts on the doublet index, on each of the 16 spinor components alike
& new & \\
\(\mathcal W_{U2A}\), \(\mathcal F_{U2AC}\) & the connection as a
Hermitian \(2\times2\) matrix,
\(\mathcal W_{U2A}=\tfrac12g_{U2}\boldsymbol\tau_{U2}\cdot\boldsymbol W_{U2A}+g'_{U2}y_{U2}B_{U2A}\mathsf 1_2\)
(\(D_{U2A}=\partial_A+i\mathcal W_{U2A}\)), and its field strength
\(\mathcal F_{U2AC}=\partial_A\mathcal W_{U2C}-\partial_C\mathcal W_{U2A}+i[\mathcal W_{U2A},\mathcal W_{U2C}]\)
& new & \texttt{cal\_w} \\
\(\Upsilon_{U2}\); \(\boldsymbol\Theta_{U2}\), \(\Theta^0_{U2}\) & a
gauge transformation
\(\Upsilon_{U2}=\exp(\tfrac i2\boldsymbol\Theta_{U2}\cdot\boldsymbol\tau_{U2})\exp(iy_{U2}\Theta^0_{U2})\),
with a triple \(\boldsymbol\Theta_{U2}(X)\) and a function
\(\Theta^0_{U2}(X)\) of the coordinates & new & \texttt{ups}
\texttt{theta} \texttt{theta0} \\
\(J^A_{U2}\), \(\boldsymbol J^A_{U2}\) & the \(U(1)_{U2}\) current
\(\bar\Psi_{U2}\gamma^A_{1671}\Psi_{U2}\) and the \(SU(2)_{U2}\)
currents
\(\bar\Psi_{U2}\gamma^A_{1671}\tfrac12\boldsymbol\tau_{U2}\Psi_{U2}\) &
new & \\
\end{longtable}

\textbf{4 Pairs, couplings, masses and Lagrangian}

\begin{longtable}[]{@{}
  >{\raggedright\arraybackslash}p{(\columnwidth - 6\tabcolsep) * \real{0.2200}}
  >{\raggedright\arraybackslash}p{(\columnwidth - 6\tabcolsep) * \real{0.5000}}
  >{\raggedright\arraybackslash}p{(\columnwidth - 6\tabcolsep) * \real{0.0800}}
  >{\raggedright\arraybackslash}p{(\columnwidth - 6\tabcolsep) * \real{0.2000}}@{}}
\toprule\noalign{}
\begin{minipage}[b]{\linewidth}\raggedright
symbol
\end{minipage} & \begin{minipage}[b]{\linewidth}\raggedright
meaning
\end{minipage} & \begin{minipage}[b]{\linewidth}\raggedright
origin
\end{minipage} & \begin{minipage}[b]{\linewidth}\raggedright
names in programs and data files
\end{minipage} \\
\midrule\noalign{}
\endhead
\bottomrule\noalign{}
\endlastfoot
\(\Phi^{+}_{U2}\), \(\Phi^{-}_{U2}\) & the spin-zero pairs
\(\psi^{1T}_{U2}C^\pm_{1671}\psi^2_{U2}=\tfrac12\Psi^T_{U2}(\mathsf E\otimes C^\pm_{1671})\Psi_{U2}\);
\(\Phi^+_{U2}\) is a Lorentz scalar whose members have spin zero and
orbital angular momentum zero in the rest frame (the singlet of total
spin zero of the task); \(\Phi^-_{U2}\), its pseudoscalar partner,
vanishes in the non-relativistic limit; both are \(SU(2)_{U2}\) singlets
of \(U(1)_{U2}\) charge 1 & new & \\
\(m_{U2}\) & the bare Dirac mass of each member of the doublet & new &
\texttt{m} \\
\(G^S_{U2}\), \(G^P_{U2}\) & the quartic couplings of
\((\bar\Psi_{U2}\Psi_{U2})^2\) and of
\(\Phi^{+\dagger}_{U2}\Phi^+_{U2}\) in the Lagrangian; positive values
attract; dimension mass\(^{-6}\) & new & \texttt{gs} \texttt{gp} \\
\(\mathcal L_{U2}\) & the Lagrangian density of the model & new & \\
\(H_{U2}\), \(N_{U2}\) & the Hamiltonian and the fermion-number operator
(charge-symmetric ordering); the ground state minimizes
\(H_{U2}-\mu_{U2}N_{U2}\) & new & \\
\(\mathcal L^{\rm MF}_{U2}\), \(H^{\rm MF}_{U2}\) & the Lagrangian and
the Hamiltonian of the mean-field approximation, in which the squares of
the fluctuations of \(\bar\Psi_{U2}\Psi_{U2}\) and \(\Phi^+_{U2}\) are
dropped & new & \\
\(\mathrm{Vol}_{U2}\) & the volume of the seven-dimensional space (a
large box with periodic boundary conditions; at the end
\(\mathrm{Vol}_{U2}\to\infty\)) & new & \\
\end{longtable}

\textbf{5 Mean-field theory and spinor solutions}

\begin{longtable}[]{@{}
  >{\raggedright\arraybackslash}p{(\columnwidth - 6\tabcolsep) * \real{0.2200}}
  >{\raggedright\arraybackslash}p{(\columnwidth - 6\tabcolsep) * \real{0.5000}}
  >{\raggedright\arraybackslash}p{(\columnwidth - 6\tabcolsep) * \real{0.0800}}
  >{\raggedright\arraybackslash}p{(\columnwidth - 6\tabcolsep) * \real{0.2000}}@{}}
\toprule\noalign{}
\begin{minipage}[b]{\linewidth}\raggedright
symbol
\end{minipage} & \begin{minipage}[b]{\linewidth}\raggedright
meaning
\end{minipage} & \begin{minipage}[b]{\linewidth}\raggedright
origin
\end{minipage} & \begin{minipage}[b]{\linewidth}\raggedright
names in programs and data files
\end{minipage} \\
\midrule\noalign{}
\endhead
\bottomrule\noalign{}
\endlastfoot
\(\mu_{U2}\) & the chemical potential conjugate to \(N_{U2}\) & new &
\texttt{mu} \\
\(M_{U2}\) & the effective mass
\(m_{U2}-2G^S_{U2}\langle\bar\Psi_{U2}\Psi_{U2}\rangle\) (M\_sym: its
value on the lowest state with \(\Delta_{U2}=0\)) & new & \texttt{M}
\texttt{mass} \texttt{M\_sym} \\
\(\Delta_{U2}\); \(\Delta_1\), \(\Delta_2\) & the pair field
\(\Delta_{U2}=G^P_{U2}\langle\Phi^+_{U2}\rangle\), complex, the order
parameter of the breaking of \(U(2)_{U2}\); \(\Delta_1\), \(\Delta_2\)
its real and imaginary parts (in the data files: its value at the
minimum; lg: \(\log_{10}\), ln: natural logarithm, Landau: the small-gap
form) & new & \texttt{Delta} \texttt{Delta\_R7} \texttt{Delta\_R8}
\texttt{Delta\_chiral\_R7} \texttt{Landau\_R7} \texttt{Landau\_R8}
\texttt{Delta1} \texttt{Delta2} \texttt{D050} \texttt{D080}
\texttt{D095} \texttt{D105} \texttt{lgD050} \texttt{lgD080}
\texttt{lgD095} \texttt{lgD105} \texttt{lnDelta} \\
\(s_{U2}\) & the sum \(M_{U2}^2+\lvert\Delta_{U2}\rvert^2\), on which
the vacuum spectrum depends & new & \texttt{s} \\
\(\vec q\), \(q\), \(q^j\) & the spatial momentum of a fermion mode
(seven components \(q^j\), \(j=0,1,2,3,5,6,7\)) and its length & new &
\texttt{q} \\
\(\psi^1_{U2}(\vec q)\), \(\psi^2_{U2}(\vec q)\) & the Fourier
components
\(\mathrm{Vol}_{U2}^{-1/2}\int d^7x\,\exp(-i\vec q\cdot\vec x)\,\psi^{1}_{U2}\)
(and the same for \(\psi^2_{U2}\)) at a fixed time \(x4\) & new & \\
\(q^F_{U2}\) & the Fermi momentum \(\sqrt{\mu_{U2}^2-M_{U2}^2}\) & new
& \\
\(\vec Q_{U2}\) & the pair momentum of a condensate
\(\Delta_{U2}\exp(i\vec Q_{U2}\cdot\vec x)\) & new & \\
\(\varepsilon_{U2}\) & the single-particle energy
\(\sqrt{q^2+M_{U2}^2}\) & new & \\
\(\mathsf h_{U2}\) & the \(16\times16\) Hamiltonian matrix
\(\alpha^j_{1671}q^j+\beta_{1671}M_{U2}\) (sum over the seven space-like
\(j\)), eigenvalues \(\pm\varepsilon_{U2}\), each 8-fold & new & \\
\(\mathsf e^{+}_{U2}\), \(\mathsf e^{-}_{U2}\) & real orthonormal
eigenvectors of \(\mathsf h_{U2}(\vec q)\) with the eigenvalues
\(+\varepsilon_{U2}\) and \(-\varepsilon_{U2}\) (eight of each sign) &
new & \\
\(\mathcal N_{U2}(\vec q)\) & the Nambu spinor
\((\psi^1_{U2}(\vec q),\ \psi^{2}_{U2}(-\vec q)^{\dagger T})\): 32
components, the annihilators of \(\psi^1_{U2}\) at \(\vec q\) and the
creators of \(\psi^2_{U2}\) at \(-\vec q\) & new & \\
\(\mathsf H_{U2}\) & the real symmetric \(32\times32\) mean-field
(Bogoliubov-de Gennes) matrix of \(\mathcal N_{U2}(\vec q)\): diagonal
blocks \(\mathsf h_{U2}-\mu_{U2}\) and
\(\alpha^j_{1671}q^j-\beta_{1671}M_{U2}+\mu_{U2}\), off-diagonal blocks
\(\Delta_{U2}C^+_{1671}\) (for real \(\Delta_{U2}\)) & new & \\
\(\mathsf Y_{U2}\), \(\mathsf Y'_{U2}\) & the \(32\times32\) matrices of
the cross terms in the square of \(\mathsf H_{U2}\), for the spin-zero
pair (\(\mathsf Y_{U2}\), square
\(\mu^2_{U2}\varepsilon^2_{U2}\mathsf 1_{32}\)) and for the pseudoscalar
pair (\(\mathsf Y'_{U2}\), square
\((\mu^2_{U2}\varepsilon^2_{U2}+M^2_{U2}\lvert\Delta_{U2}\rvert^2)\mathsf 1_{32}\))
& new & \\
\(E^{+}_{U2}\), \(E^{-}_{U2}\) & the quasi-particle energies
\(\sqrt{(\varepsilon_{U2}\pm\mu_{U2})^2+\lvert\Delta_{U2}\rvert^2}\)
(\(E^+_{U2}\): the branch of the antiparticles, \(E^-_{U2}\): the branch
of the particles; E\_Cminus: the smaller energy of the pseudoscalar
pairing) & new & \texttt{E\_Cplus} \texttt{E\_Cminus} \\
\(\xi^{\pm}_{U2}\), \(u^{\pm}_{U2}\), \(v^{\pm}_{U2}\) &
\(\xi^\pm_{U2}=\pm\varepsilon_{U2}-\mu_{U2}\) and the Bogoliubov
coefficients \(u^\pm_{U2}=[(1+\xi^\pm_{U2}/E^\mp_{U2})/2]^{1/2}\),
\(v^\pm_{U2}=[(1-\xi^\pm_{U2}/E^\mp_{U2})/2]^{1/2}\) (note
\(E^\mp_{U2}=[(\xi^\pm_{U2})^2+\lvert\Delta_{U2}\rvert^2]^{1/2}\)) & new
& \\
\(a_{U2}\), \(d_{U2}\) & for an eigenvector \(\mathsf e^\pm_{U2}\) of
\(\mathsf h_{U2}(\vec q)\): the mode annihilators
\(a_{U2}=(\mathsf e^\pm_{U2})^T\psi^1_{U2}(\vec q)\) and
\(d_{U2}=(\beta_{1671}\mathsf e^\pm_{U2})^T\psi^2_{U2}(-\vec q)\) & new
& \\
\(\lvert\mathrm{SF}_{U2}\rangle\), \(b^\dagger_{U2}\),
\(\lvert\varnothing_{U2}\rangle\) & the superfluid ground state; the
operator \(b^\dagger_{U2}=a^\dagger_{U2}d^\dagger_{U2}\) that creates
one spin-zero pair of members of opposite momenta; the state without any
fermion in any mode (annihilated by every \(a_{U2}\) and \(d_{U2}\));
the ground state is a product of factors
\(u^\pm_{U2}-v^\pm_{U2}b^\dagger_{U2}\) applied to it & new & \\
\end{longtable}

\textbf{6 Effective potential and its minima}

\begin{longtable}[]{@{}
  >{\raggedright\arraybackslash}p{(\columnwidth - 6\tabcolsep) * \real{0.2200}}
  >{\raggedright\arraybackslash}p{(\columnwidth - 6\tabcolsep) * \real{0.5000}}
  >{\raggedright\arraybackslash}p{(\columnwidth - 6\tabcolsep) * \real{0.0800}}
  >{\raggedright\arraybackslash}p{(\columnwidth - 6\tabcolsep) * \real{0.2000}}@{}}
\toprule\noalign{}
\begin{minipage}[b]{\linewidth}\raggedright
symbol
\end{minipage} & \begin{minipage}[b]{\linewidth}\raggedright
meaning
\end{minipage} & \begin{minipage}[b]{\linewidth}\raggedright
origin
\end{minipage} & \begin{minipage}[b]{\linewidth}\raggedright
names in programs and data files
\end{minipage} \\
\midrule\noalign{}
\endhead
\bottomrule\noalign{}
\endlastfoot
\(\Omega_{U2}\); \(\Omega^{\rm cl}_{U2}\) & the effective potential of
the superfluid: the grand potential per unit volume at zero temperature
as a function of \(M_{U2}\) and \(\Delta_{U2}\) (data files: relative to
the normal state of mass \(m_{U2}\) at the same \(\mu_{U2}\); lNNN: at
\(\lambda^P_{U2}=\) N.NN);
\(\Omega^{\rm cl}_{U2}=(M_{U2}-m_{U2})^2/(4G^S_{U2})+\lvert\Delta_{U2}\rvert^2/G^P_{U2}\)
its classical part & new & \texttt{Omega} \texttt{Omega\_R7}
\texttt{Omega\_R8} \texttt{O050} \texttt{O080} \texttt{O095}
\texttt{O105} \texttt{l095} \texttt{l100} \texttt{l102} \texttt{l105}
\texttt{l110} \\
\(\Omega^{\min}_{U2}\), \(\Omega^{\rm sym}_{U2}\) & the value at the
global minimum, and on the lowest state with \(\Delta_{U2}=0\)
(\(U(2)_{U2}\) unbroken) & new & \texttt{Omega\_sym} \\
\(\Lambda_{U2}\) & the momentum cutoff of the model; every number of
Revision\_3 is in units of it & new & \\
\(G^{\rm crit}_{U2}\) & the critical pair coupling of the vacuum:
\(90\pi^4\Lambda_{U2}^{-6}\) (R7), \(288\pi^4\Lambda_{U2}^{-6}\) (R8);
twice the value for the left-handed doublet & new & \texttt{Gcrit\_R7}
\texttt{Gcrit\_R8} \texttt{Gcrit\_chiral\_R7} \\
\(\lambda^P_{U2}\), \(\lambda^S_{U2}\) & the dimensionless couplings
\(G^P_{U2}/G^{\rm crit}_{U2}\) and \(4G^S_{U2}/G^{\rm crit}_{U2}\)
(inv\_G: \(1/G^P_{U2}\)) & new & \texttt{lambda} \texttt{inv\_G} \\
\(a^{(7)}_{U2}\), \(a^{(8)}_{U2}\) & \(A_7/(2\pi)^7=1/(120\pi^4)\) and
\(A_8/(2\pi)^8=1/(768\pi^4)\): for an integrand that depends on the
length \(q\) only,
\(\int d^7q/(2\pi)^7(\dots)=a^{(7)}_{U2}\int_0^{\Lambda_{U2}}q^6(\dots)\,dq\),
and the same in 8 dimensions with \(\kappa_{U2}\) & new & \texttt{A7} \\
\(\kappa_{U2}\), \(\kappa^{(4)}_{U2}\) & the length of the Euclidean
8-momentum, whose components are the seven components \(q^j\) and the
imaginary-time frequency \(\kappa^{(4)}_{U2}\), in the Lorentz-invariant
regularization (R8); \(\kappa^2_{U2}=q^2+(\kappa^{(4)}_{U2})^2\) & new
& \\
\(L_{U2}\) & the function
\(L_{U2}(s_{U2})=2\int_0^1\kappa_{U2}^7\ln(\kappa_{U2}^2+s_{U2})\,d\kappa_{U2}\)
of the vacuum potential of regularization (R8) (units
\(\Lambda_{U2}=1\)) & new & \\
\(\mathcal J^{(\iota)}_{U2}\); \(\iota\) & the radial integrals
\(\mathcal J^{(\iota)}_{U2}(s_{U2})=\int_0^1q^\iota(q^2+s_{U2})^{-1/2}dq\)
of the vacuum potential of regularization (R7), for the even exponents
\(\iota=0,2,4,6,8\) & new & \\
\(n_{U2}\) & the number density
\(-\partial\Omega_{U2}/\partial\mu_{U2}\) of the fermions
(charge-symmetric: zero in the vacuum) & new & \texttt{n} \texttt{n050}
\texttt{n080} \texttt{n095} \texttt{n105} \texttt{n\_over\_4mu} \\
\(z_{U2}\) & \(\Delta_{U2}\,d\ln\Delta_{U2}/d\mu_{U2}\) along the
minimum: how fast the condensation energy changes with \(\mu_{U2}\) &
new & \texttt{dlnDelta\_dmu} \texttt{Delta\_dlnDelta\_dmu} \\
\end{longtable}

\textbf{7 Collective modes, stiffness and boson masses}

\begin{longtable}[]{@{}
  >{\raggedright\arraybackslash}p{(\columnwidth - 6\tabcolsep) * \real{0.2200}}
  >{\raggedright\arraybackslash}p{(\columnwidth - 6\tabcolsep) * \real{0.5000}}
  >{\raggedright\arraybackslash}p{(\columnwidth - 6\tabcolsep) * \real{0.0800}}
  >{\raggedright\arraybackslash}p{(\columnwidth - 6\tabcolsep) * \real{0.2000}}@{}}
\toprule\noalign{}
\begin{minipage}[b]{\linewidth}\raggedright
symbol
\end{minipage} & \begin{minipage}[b]{\linewidth}\raggedright
meaning
\end{minipage} & \begin{minipage}[b]{\linewidth}\raggedright
origin
\end{minipage} & \begin{minipage}[b]{\linewidth}\raggedright
names in programs and data files
\end{minipage} \\
\midrule\noalign{}
\endhead
\bottomrule\noalign{}
\endlastfoot
\(\sigma_{U2}\), \(m_{\sigma U2}\) & the amplitude mode of the pair
field about the minimum: a real scalar boson, neutral under every gauge
generator, that exists only after the breaking; its mass (threshold:
twice the smallest quasi-particle energy; kind 0: below it, 1: at it) &
new & \texttt{msigma} \texttt{msigma\_over\_2Delta} \texttt{threshold}
\texttt{kind} \\
\(\theta_{U2}\) & the phase of the pair field; its fluctuation is the
Goldstone mode, which \(B_{U2A}\) absorbs & new & \\
\(f^s_{U2}\), \(f^t_{U2}\), \(f_{U2}\) & the stiffness of the phase: the
coefficients of \(\tfrac12\sum_j(\partial_j\theta_{U2})^2\) and
\(\tfrac12(\partial_4\theta_{U2})^2\) in the effective Lagrangian; in
the Lorentz-invariant regularization (R8) both equal \(f_{U2}\) (data
files: their squares) & new & \texttt{fs2} \texttt{ft2}
\texttt{ft2\_fixed} \texttt{ft2\_relaxed} \texttt{f2\_R8}
\texttt{fs2\_over\_n4mu} \\
\(\mathcal L^{\rm eff}_{U2}\) & the effective Lagrangian of the phase
\(\theta_{U2}\) and of \(B_{U2A}\) at long wavelengths & new & \\
\(m_{BU2}\) & the mass \(g'_{U2}f_{U2}\) of \(B_{U2A}\) after the
breaking (the pair has the charge 1) & new & \\
\(v^G_{U2}\) & the velocity \(f^s_{U2}/f^t_{U2}\) of the Goldstone mode
if \(U(1)_{U2}\) were not gauged (vG2: its square) & new &
\texttt{vG\_R7} \texttt{vG2} \\
\(w_{U2}\) & the frequency of a small oscillation of the composite
fields (time dependence \(\exp(-iw_{U2}x4)\)) & new & \\
\(\mathsf V^1_{U2}\), \(\mathsf V^2_{U2}\), \(\mathsf V^3_{U2}\) & the
\(32\times32\) vertices of \(\Delta_1\), \(\Delta_2\) and \(M_{U2}\) in
\(\mathsf H_{U2}\):
\(\mathsf V^1_{U2}=\begin{pmatrix}0&C^+_{1671}\\C^+_{1671}&0\end{pmatrix}\),
\(\mathsf V^2_{U2}=i\begin{pmatrix}0&C^+_{1671}\\-C^+_{1671}&0\end{pmatrix}\),
\(\mathsf V^3_{U2}=\begin{pmatrix}\beta_{1671}&0\\0&-\beta_{1671}\end{pmatrix}\)
& new & \\
\(\lvert\mathrm o\rangle_{U2}\), \(\lvert\mathrm e\rangle_{U2}\);
\(\mathcal E^{\rm o}_{U2}\), \(\mathcal E^{\rm e}_{U2}\) & an occupied
and an empty eigenvector of \(\mathsf H_{U2}(\vec q)\) and their
eigenvalues, \(\mathcal E^{\rm o}_{U2}<0\le\mathcal E^{\rm e}_{U2}\) &
new & \\
\(\Pi_{U2}\), \(\mathsf R_{U2}\) & the \(3\times3\) polarization matrix
of the fermions for the composite fields \((\Delta_1,\Delta_2,M_{U2})\)
(entries numbered 1, 2, 3 in this order) and the matrix
\(\mathsf R_{U2}(w_{U2})\) of the small oscillations, the second
derivatives of \(\Omega^{\rm cl}_{U2}\) plus \(\Pi_{U2}(w_{U2})\)
(random-phase approximation) & new & \\
\(c^\sigma_{U2}\) & the positive function
\(c^\sigma_{U2}(w_{U2})=16\int\frac{d^7q}{(2\pi)^7}(q^2+s_{U2})^{-1/2}\big[4(q^2+s_{U2})-w_{U2}^2\big]^{-1}\),
with
\([\mathsf R_{U2}(w_{U2})]_{11}=c^\sigma_{U2}(w_{U2})(4\lvert\Delta_{U2}\rvert^2-w_{U2}^2)\)
in the vacuum without scalar coupling & new & \\
\(\mathcal V^{\sigma}_{U2}\) & the pair interaction carried by
\(\sigma_{U2}\), \(-[\mathsf R_{U2}(w_{U2})^{-1}]_{11}\); at
\(w_{U2}=0\) it is \(-[\partial^2\Omega_{U2}/\partial\Delta_1^2]^{-1}\)
& new & \\
\(V^{\rm C}_{U2}\), \(\phi^{\rm Y}_{U2}\), \(V^{\rm pair}_{U2}\) & the
potential
\(V^{\rm C}_{U2}(\varrho_{U2})=\Gamma_{\rm E}(5/2)/(4\pi^{7/2}\varrho_{U2}^5)\)
of a massless vector boson between two unit charges in seven space
dimensions; the factor
\(\phi^{\rm Y}_{U2}(m_{BU2}\varrho_{U2})=\exp(-m_{BU2}\varrho_{U2})\)
times \([(m_{BU2}\varrho_{U2})^2+3m_{BU2}\varrho_{U2}+3]/3\) by which
the mass of \(B_{U2A}\) multiplies it; the static potential of two heavy
members of the doublet (data files: x is \(m_{BU2}\varrho_{U2}\)) & new
& \texttt{x} \texttt{phi} \texttt{singlet\_after}
\texttt{triplet\_after} \texttt{singlet\_before}
\texttt{triplet\_before} \\
\end{longtable}

\textbf{8 Numbers, regularizations and labels}

\begin{longtable}[]{@{}
  >{\raggedright\arraybackslash}p{(\columnwidth - 6\tabcolsep) * \real{0.2200}}
  >{\raggedright\arraybackslash}p{(\columnwidth - 6\tabcolsep) * \real{0.5000}}
  >{\raggedright\arraybackslash}p{(\columnwidth - 6\tabcolsep) * \real{0.0800}}
  >{\raggedright\arraybackslash}p{(\columnwidth - 6\tabcolsep) * \real{0.2000}}@{}}
\toprule\noalign{}
\begin{minipage}[b]{\linewidth}\raggedright
symbol
\end{minipage} & \begin{minipage}[b]{\linewidth}\raggedright
meaning
\end{minipage} & \begin{minipage}[b]{\linewidth}\raggedright
origin
\end{minipage} & \begin{minipage}[b]{\linewidth}\raggedright
names in programs and data files
\end{minipage} \\
\midrule\noalign{}
\endhead
\bottomrule\noalign{}
\endlastfoot
(R7), (R8) & the two regularizations: (R7) a sharp cutoff on the length
of the seven-dimensional spatial momentum; (R8) a sharp cutoff on the
length of the Euclidean 8-momentum (Lorentz invariant; used in the
vacuum) & new & \texttt{R7} \texttt{R8} \\
\(\Gamma_{\rm E}\) & Euler's Gamma function;
\(\Gamma_{\rm E}(5/2)=3\sqrt\pi/4\) & new & \\
\(i\) & the imaginary unit & new & \\
\({}^T\), \({}^\dagger\), \({}^*\); \(\otimes\); \([\ ,\ ]\),
\(\{\ ,\ \}\); \(\langle\ \rangle\) & transpose, Hermitian conjugate,
complex conjugate (for operators: the Hermitian conjugate of each
component); Kronecker product; commutator and anticommutator;
expectation value in the ground state & new & \\
\end{longtable}

\textbf{9 The electroweak theory as the template (suffix SM)}

\begin{longtable}[]{@{}
  >{\raggedright\arraybackslash}p{(\columnwidth - 6\tabcolsep) * \real{0.2200}}
  >{\raggedright\arraybackslash}p{(\columnwidth - 6\tabcolsep) * \real{0.5000}}
  >{\raggedright\arraybackslash}p{(\columnwidth - 6\tabcolsep) * \real{0.0800}}
  >{\raggedright\arraybackslash}p{(\columnwidth - 6\tabcolsep) * \real{0.2000}}@{}}
\toprule\noalign{}
\begin{minipage}[b]{\linewidth}\raggedright
symbol
\end{minipage} & \begin{minipage}[b]{\linewidth}\raggedright
meaning
\end{minipage} & \begin{minipage}[b]{\linewidth}\raggedright
origin
\end{minipage} & \begin{minipage}[b]{\linewidth}\raggedright
names in programs and data files
\end{minipage} \\
\midrule\noalign{}
\endhead
\bottomrule\noalign{}
\endlastfoot
\(\phi_{\rm SM}\), \(Y_{\rm SM}\) & the Higgs doublet of the electroweak
theory and its hypercharge \(Y_{\rm SM}=\tfrac12\) (quoted as the
template of the new covariant derivative) & new & \\
\(\mathrm D^{\rm SM}\), \(g_{\rm SM}\), \(g'_{\rm SM}\),
\(\boldsymbol\tau_{\rm SM}\), \(\boldsymbol W_{\rm SM}\), \(B_{\rm SM}\)
& the covariant exterior derivative of the electroweak theory, its two
couplings, the Pauli matrices and the gauge fields (one-forms),
\(\mathrm D^{\rm SM}\phi_{\rm SM}=(\mathrm d+\tfrac i2g_{\rm SM}\boldsymbol\tau_{\rm SM}\cdot\boldsymbol W_{\rm SM}+ig'_{\rm SM}Y_{\rm SM}B_{\rm SM})\phi_{\rm SM}\);
\(\mathrm d\) is the exterior derivative & new & \\
\end{longtable}

\end{document}
